\makeabstract
    {
    Characterization of Cam Kinase II Mutant Overexpression in Drosophila NMJ
    }
    {
    Lauren Ro\textsuperscript{1}, Madison Vant\textsuperscript{1}, John P. Roche\textsuperscript{1}
    }
    {
    \textsuperscript{1} Neuroscience Program, Amherst College, Amherst, MA
    }
    {
    The neuromuscular junctions (NMJ) of Drosophila share many similarities to humans and can be used as a model for studying neurological diseases. This study focuses on the impact of Calcium/Calmodulin-Stimulated Protein Kinase II (CaMKII) on synaptic development and function; laying the foundation for future work expanding on research conducted by UMass Amherst on the ability to reverse the effects of overexpression of CaMKII.
    }
    {
    Three strains of driver flies were used: 1. OK371-Gal4 activates gene expression in glutamatergic neurons, 2. Elav-Gal4 is pan neuronal, activating gene expression in all post-mitotic neurons, or 3. C380 (X)-Gal4 activates gene expression in motor and some peripheral neurons. Each genetic cross contained three virgin females from one of the drivers and three males from either the control (CS), T287D mutant, or D136N, T287D mutant. In total, there were seven different crosses. Larvae from each of these crosses were dissected to expose the NMJ. Through immunohistochemistry using mouse \ensuremath{\alpha}-BRP and rabbit \ensuremath{\alpha}-GluRC as primary antibodies, the larvae were stained with fluorescent markers. The larvae were mounted onto a slide and imaged using a confocal microscope. The images were analyzed using Fiji to count the number of boutons, integrated density of BRP, integrated density of GluRC, and total area of the synapse. Ten larvae from each cross were used in a behavioral assay. Larvae were placed onto an agarose dish with graph paper underneath, and two trials of two minutes were performed to measure the total distance traveled. ANOVAs were used for statistical analysis using Prism.
    }
    {
    Flies with the T287D mutation inserted into both chromosome (Chr) II and III resulted in increased number of boutons at the NMJ. This occurred when using the OK371-Gal4 and elav-Gal4 driver. Additionally, larvae with the double point mutation (OK371-Gal4 x D136N, T287D) inserted into Chr III, had significantly less boutons than the T287D mutants while having no significant difference between the controls. The T287D mutation inserted into Chr II increased eclosion while when inserted into Chr III, decreased eclosion. There were no significant differences in the areas of the NMJs or locomotor behavior in larvae between crosses.
    }
    {
    Proper expression of CaMKII is essential for NMJ development and function. The number of boutons increases as a result of the overexpression of CaMKII due to the insertion of UAS-T287D CaMKII. This indicates the increase in the number of boutons is a phosphorylation dependent mechanism rather than a localization or structural mechanism. In addition, elav-Gal4 which is a pan neuronal driver, impacts synaptic development (in terms of number of boutons) to a greater extent than more selective drivers (OK371-Gal4 and C380 (X)-Gal4).
    }

    \newpage
    